% Template for Elsevier CRC journal article
% Procedia Computer Science - ICMLDE 2026
% Concise Conference Version (Target <= 10 Pages)

\documentclass[3p,times,procedia]{elsarticle}
\flushbottom

\usepackage{ecrc}
\usepackage[bookmarks=false]{hyperref}
\hypersetup{colorlinks,
  linkcolor=blue,
  citecolor=blue,
  urlcolor=blue}
\usepackage{amssymb}
\usepackage{amsmath}
\usepackage{booktabs}
\usepackage{array}
\usepackage{graphicx}

%% ecrc package configurations
\volume{00}
\firstpage{1}
\journalname{Procedia Computer Science}
\runauth{Anonymous Author(s)}
\jid{procs}

%% Vancouver numbered citations
\biboptions{sort&compress}

\begin{document}
\begin{frontmatter}

\dochead{5th International Conference on Machine Learning and Data Engineering}

\title{Source-Stratified Evaluation of Out-of-Distribution Detection in Fruit-Quality Classification Under Species and Treatment-Labelled Shifts}

\author[a]{Author Names Redacted for Peer Review}
\address[a]{Affiliation Details Redacted for Double-Blind Peer Review}

\begin{abstract}
Automated fruit-quality classification via computer vision achieves high accuracy on curated benchmarks but frequently degrades under real-world distribution shifts. This study conducts a rigorous, source-stratified empirical evaluation of a frozen MobileNetV2 six-class fruit-quality classifier and five post-hoc out-of-distribution (OOD) scoring methods: Maximum Softmax Probability (MSP), Energy, an ODIN-style perturbation variant, Ledoit--Wolf Mahalanobis distance, and a modified OpenMax-style heuristic, alongside an operational Energy--MSP Boolean Dual-Gate cascade. The evaluation encompasses $N = 13,465$ photographic records across three repositories, organized into primary in-distribution, external known-class, unseen-species, dataset-reported treatment-labelled, and combined-shift cohorts. Closed-set accuracy degrades from 96.92\% on primary laboratory test data to 68.36\% on pooled external known-class imagery (a 28.56 percentage-point decline), driven primarily by an asymmetric collapse on Fresh Apples (32.12\% accuracy; 310/965 correct). For semantic species shift within FruitVision (Cohort A), Energy achieves 81.57\% AUROC, whereas on Sultana (Cohort D) it drops to 67.52\%. Under dataset-reported chemical treatment shift (Cohort B), all evaluated post-hoc scoring functions exhibit weak discrimination (Energy AUROC = 53.13\%, Mahalanobis = 47.91\%). Under validation-calibrated operating thresholds ($\tau_{95}$), the Dual-Gate incurs a 42.43\% known-class rejection rate on external imagery and accepts 42.83\% (821/1,917) of treatment-labelled image records into classifier Fresh-class output labels. Perceptual-hash sensitivity analysis establishes detector ranking invariance across deduplication tiers (maximum AUROC shift $\le 0.22$ percentage points). These findings expose the visual failure boundaries of standalone RGB classification and demonstrate the value of source-stratified evaluation; future work could evaluate whether complementary sensing modalities improve screening in post-harvest grading.
\end{abstract}

\begin{keyword}
Fruit quality classification \sep Open-set recognition \sep Out-of-distribution detection \sep Domain shift \sep MobileNetV2 \sep Energy-based detection \sep Food quality assessment
\end{keyword}

\cortext[cor1]{Corresponding author. Contact details redacted for double-blind peer review.}

\end{frontmatter}

\section{Introduction}
\label{sec:introduction}

Post-harvest agricultural losses present a major challenge to global food security and supply chain sustainability \cite{fao2019state}. The Food and Agriculture Organization (FAO) estimates that approximately 14\% of global food production is lost between harvest and retail distribution \cite{fao2019state}, with commercial fruit losses consistently ranging between 7\% and 16\% (e.g., 10.39\% in apple, 7.76\% in banana, and 15.88\% in citrus) \cite{jha2015assessment}. Manual visual sorting is labor-intensive, subjective, and prone to inspector fatigue \cite{rizzo2023fruit,zhang2024fruit}. While convolutional neural networks (CNNs) such as MobileNetV2 \cite{sandler2018mobilenetv2} have achieved over 95\% accuracy on curated fruit-quality benchmarks \cite{chen2024classification,gulzar2023fruit}, standard softmax classifiers suffer from severe overconfidence: normalized probability distributions peak sharply near 1.0 even on novel species or corrupted imagery \cite{guo2017calibration,hendrycks2017baseline}. In automated packing facilities, an overconfident false positive---such as misrouting an unseen species or a chemically adulterated lot into a fresh output stream---creates severe commercial and food-quality liabilities.

Post-hoc out-of-distribution (OOD) detection algorithms aim to identify and reject anomalous inputs at inference time without requiring model retraining \cite{scheirer2013toward,bendale2016towards,yang2024generalized}. Methods such as Maximum Softmax Probability (MSP) \cite{hendrycks2017baseline}, input-perturbed temperature scaling (ODIN) \cite{liang2018enhancing}, Mahalanobis feature distance \cite{lee2018simple}, and Energy-based detection \cite{liu2020energy} have demonstrated success on synthetic benchmarks \cite{vaze2022open}. However, existing agricultural evaluations conflate distribution shifts: when anomalies are drawn from external repositories while in-distribution data originates from curated laboratory archives, detector performance is confounded by differences in camera hardware, ambient illumination, and background environments rather than true botanical or pathological divergence.

To resolve these limitations, this study conducts a source-stratified empirical benchmark evaluating a frozen MobileNetV2 six-class fruit-quality classifier across $N = 13,465$ photographic records from three independent repositories: Kaggle Fruit Quality, FruitVision \cite{bijoy2025fruitvision}, and Sultana et al.\ \cite{sultana2022extensive}. By enforcing source-matched comparisons, we isolate external domain shift from semantic species and dataset-reported chemical treatment shifts, addressing three core research questions:
\begin{itemize}
\setlength{\itemsep}{0pt}\setlength{\parsep}{0pt}\setlength{\parskip}{0pt}\setlength{\topsep}{1pt}\setlength{\partopsep}{0pt}
\item \textbf{RQ1 (External Domain Robustness):} How severely does closed-set accuracy degrade when transferring from laboratory benchmarks to external photographic repositories, and what are the specific class-wise failure modes?
\item \textbf{RQ2 (Comparative OOD Discrimination):} How effectively do post-hoc scoring functions separate in-distribution produce from unseen botanical species versus dataset-reported chemical preservative shifts under source-matched conditions?
\item \textbf{RQ3 (Operational Gating Feasibility and Trade-offs):} What operational rejection and false-acceptance trade-offs emerge under fixed, validation-calibrated operating thresholds?
\end{itemize}

\noindent\textbf{Contributions:} (1) A source-stratified evaluation framework across five partitions and four cohorts spanning $13,465$ images. (2) Diagnostic quantification of a 28.56 percentage-point closed-set accuracy drop on external data, driven by an asymmetric collapse on Fresh Apples (32.12\% accuracy). (3) Evaluation of five post-hoc OOD algorithms and an operational Energy--MSP Dual-Gate cascade, demonstrating moderate discrimination on semantic species shifts (AUROC up to 81.57\%) but failure on dataset-reported chemical treatment shifts (AUROC $\le 58.77\%$, Dual-Gate UFAR-Fresh $= 42.83\%$). (4) Perceptual-hash sensitivity analysis establishing that detector rankings remain strictly invariant across deduplication tiers.

\section{Related Work}
\label{sec:related_work}

\noindent\textbf{Post-Harvest Vision and Post-Hoc OOD Detection:} Deep convolutional architectures have advanced automated produce sorting, defect identification, and ripeness estimation \cite{rizzo2023fruit,zhang2024fruit,chen2024classification,gulzar2023fruit}. Lightweight convolutional backbones such as MobileNetV2 \cite{sandler2018mobilenetv2} are frequently explored for embedded sorting applications \cite{chen2024classification,gulzar2023fruit} due to compact parameter footprints (approximately $2.26\times 10^6$ total parameters in our six-class configuration) and low computational overhead. However, closed-set classifiers cannot detect anomalous or non-target produce \cite{guo2017calibration}. Post-hoc OOD detection algorithms compute an anomaly score $S(x)$ from internal activations without retraining \cite{yang2024generalized}. Baseline Maximum Softmax Probability (MSP) \cite{hendrycks2017baseline} evaluates $\max_k p_k(x)$. Liang et al.\ \cite{liang2018enhancing} introduced ODIN, combining temperature scaling with input gradient perturbations: $\tilde{x} = x - \varepsilon \operatorname{sign}(-\nabla_x \log p_{\hat{y}}(x; T))$. Bendale and Boult \cite{bendale2016towards} formulated OpenMax by fitting Weibull distributions to penultimate features. Lee et al.\ \cite{lee2018simple} proposed class-conditional Gaussian modeling via Mahalanobis distance $d_M^2(x) = \min_k (f(x) - \mu_k)^\top \mathbf{\Sigma}^{-1} (f(x) - \mu_k)$. Liu et al.\ \cite{liu2020energy} proposed an energy-based OOD scoring function $S_{\text{Energy}}(x) = T \log \sum_{k=1}^K \exp(z_k(x)/T)$ computed over raw logits. The Energy score is a post-hoc logit-based OOD scoring function motivated by an energy-based formulation; it is not treated here as a calibrated estimate of the true input marginal density. Negated Euclidean Mahalanobis distance $-\sqrt{d_M^2(x)}$ strictly preserves sample ranking and AUROC while scaling linearly with feature deviations.

\vspace{2pt}
\noindent\textbf{Agricultural Domain Shifts and Research Positioning:} Recent studies explore domain adaptation and anomaly rejection in agricultural contexts. Wang et al.\ \cite{wang2023crossdomain} evaluated attention networks with unsupervised domain adaptation for fruit classification across lighting conditions. Pehlivan and Kara \cite{pehlivan2026lightweight} proposed one-shot open-set metric learning for food recognition. Jansi and Kavitha \cite{jansi2026peos} developed PEOS-Net, combining prototype learning and energy scores for crop disease recognition, while Yan et al.\ \cite{yan2026pagmsbnet} introduced PAG-MSBNet for open-set crop disease diagnosis. Despite these advances, many evaluations assess models on curated test splits collected under uniform protocols, which can leave camera hardware and background differences entangled with botanical distribution shifts. In contrast, this study evaluates a frozen standard MobileNetV2 classifier across multiple external repositories, establishing a source-stratified benchmark that systematically contrasts external covariate shift, semantic species shift, and dataset-reported chemical treatment shift under prospective, locked operating thresholds.

\section{Dataset Architecture and Experimental Design}
\label{sec:dataset_and_design}

\subsection{Photographic Repositories and Deduplication Auditing}
Our empirical benchmark evaluates $N = 13,465$ authentic photographic records from three independent sources:
\begin{itemize}
\setlength{\itemsep}{0pt}\setlength{\parsep}{0pt}\setlength{\parskip}{0pt}\setlength{\topsep}{1pt}\setlength{\partopsep}{0pt}
\item \textbf{Kaggle Fruit Quality Benchmark:} Curated benchmark covering three commercial fruit species---apple (\textit{Malus domestica}), banana (\textit{Musa acuminata}), and orange (\textit{Citrus sinensis})---in fresh and rotten states under neutral backgrounds. From 2,043 test images, we isolate strictly the 227 raw photographic images (`\texttt{is\_raw == 1}') to construct Partition A (Primary ID Test), setting aside 1,816 synthetic augmentations. An independent validation partition of 227 raw images is maintained strictly for prospective threshold calibration.
\item \textbf{FruitVision Dataset \cite{bijoy2025fruitvision}:} Large-scale benchmark containing 10,154 published raw photographs, yielding $N = 10,045$ evaluated records after 109 hash-identified duplicate exclusions across five species: three known species (apple, banana, orange) and two unseen non-target species (grape \textit{Vitis vinifera} and mango \textit{Mangifera indica}). Photographed across wholesale markets, retail stalls, and laboratory benches using three commercial smartphones under unconstrained lighting. Crucially, FruitVision includes produce lots subjected to post-harvest preservative immersion (formalin chemical treatment protocol), providing labelled condition shifts.
\item \textbf{Sultana et al.\ Dataset \cite{sultana2022extensive}:} Benchmark containing approximately 3,200 raw photographs, yielding $N = 3,193$ evaluated records after seven hash-identified duplicate exclusions of fresh and rotten fruits collected across local market stalls and open fields in Bangladesh, providing $N = 1,193$ known-class instances (apple, banana, orange) and $N = 2,000$ non-target agricultural commodities (guava, jujube, pomegranate, etc.).
\end{itemize}
MD5 checksums confirmed zero bit-level collisions across partitions. Discrete cosine transform perceptual hashing (64-bit pHash) audited image similarity across Hamming scan radii $r \in \{0, 4, 8\}$. The canonical evaluation tiers evaluate metric sensitivity: Primary ($N = 13,465$); Tier 1 ($N = 13,409$), which excludes 56 unique cross-partition collision records detected at $r \le 8$ (11 FruitVision Known reference records, 14 FruitVision D1 records, 11 FruitVision C records, 20 FruitVision D2 records; zero corresponding Sultana exclusions); and Tier 2 ($N = 13,382$), which excludes an additional 27 intra-partition multi-label collision records (6 Partition B records, 5 Partition C records, 16 Partition D1 records).

\begin{table*}[!t]
\caption{Dataset partitioning architecture, source provenance, and source-stratified evaluation cohort definitions ($N = 13,465$ photographic records).}
\label{tab:dataset_partitions_cohorts}
\centering
\footnotesize
\setlength{\tabcolsep}{3.5pt}
\resizebox{\columnwidth}{!}{%
\begin{tabular}{lllccl}
\toprule
\textbf{Partition / Cohort} & \textbf{Upstream Repository} & \textbf{Species Included} & \textbf{Condition / State} & \textbf{$N$} & \textbf{Evaluation Role / Comparison} \\
\midrule
\textbf{Partition A} & Kaggle Benchmark & Apple, Banana, Orange & Fresh, Rotten (Lab) & 227 & Primary In-Distribution (ID) Test Split \\
\textbf{Partition B (FruitVision B)} & FruitVision \cite{bijoy2025fruitvision} & Apple, Banana, Orange & Fresh, Rotten (Natural) & 4,157 & Source-matched ID reference (Cohorts A, B, C) \\
\textbf{Partition B (Sultana B)} & Sultana et al.\ \cite{sultana2022extensive} & Apple, Banana, Orange & Fresh, Rotten (Market/Field) & 1,193 & Source-matched ID reference (Cohort D) \\
\textbf{Partition C} & FruitVision \cite{bijoy2025fruitvision} & Apple, Banana, Orange & Formalin-Treated Knowns & 1,917 & Condition Shift OOD (Cohort B) \\
\textbf{Partition D1 (FruitVision)} & FruitVision \cite{bijoy2025fruitvision} & Grape, Mango & Natural Unseen Species & 2,776 & Semantic Species Shift OOD (Cohort A) \\
\textbf{Partition D1 (Sultana)} & Sultana et al.\ \cite{sultana2022extensive} & Guava, Jujube, etc. & Unseen Agricultural Produce & 2,000 & Semantic Species Shift OOD (Cohort D) \\
\textbf{Partition D2} & FruitVision \cite{bijoy2025fruitvision} & Grape, Mango & Formalin-Treated Unseen & 1,195 & Combined Semantic + Condition Shift (Cohort C) \\
\midrule
\textbf{Cohort A: Semantic} & FruitVision (Same Source) & ID: Apple, Banana, Orange & OOD: Grape, Mango (Natural) & 6,933 & Isolates pure botanical species shift \\
\textbf{Cohort B: Condition} & FruitVision (Same Source) & ID: Natural Knowns & OOD: Formalin-Treated Knowns & 6,074 & Evaluates post-harvest treatment shift \\
\textbf{Cohort C: Combined} & FruitVision (Same Source) & ID: Natural Knowns & OOD: Formalin-Treated Unseen & 5,352 & Evaluates joint species and treatment shift \\
\textbf{Cohort D: Semantic} & Sultana et al.\ (Same Source) & ID: Apple, Banana, Orange & OOD: Non-target Produce & 3,193 & Sultana source-matched semantic-shift evaluation \\
\bottomrule
\end{tabular}%
}
\end{table*}

\subsection{Partitioning Architecture and Source-Matched Cohorts}
To prevent camera hardware and lighting differences from confounding anomaly detection, we organize image records into five partitions and four source-matched cohorts (Table~\ref{tab:dataset_partitions_cohorts}). Source matching ensures that known and anomaly distributions share the same capture repository, reducing repository-level acquisition divergence.

\noindent\textbf{Specimen Independence Caveat:} Public crowdsourced repositories lack individual fruit tracking identifiers; multiple records may represent rotated captures of the same fruit. All bootstrap confidence intervals are strictly conditional on image-level resampling exchangeability.

\section{Methodology}
\label{sec:methodology}

\subsection{Classifier Architecture and Preprocessing Pipeline}
We evaluate a frozen MobileNetV2 architecture \cite{sandler2018mobilenetv2} trained on curated produce imagery. Input images are resized to $224 \times 224 \times 3$, normalized to $[-1, 1]$, and processed through inverted residual bottleneck blocks. Global average pooling produces a 1280-dimensional feature vector $f(x) \in \mathbb{R}^{1280}$, passed through a dropout layer ($p = 0.3$) and a dense softmax layer over $K = 6$ classes (Fresh/Rotten Apple, Banana, Orange). Model weights remain frozen under verified SHA-256 checksum `\texttt{4a12f354...b6b4}`.

\subsection{Evaluated Post-Hoc OOD Detection Scoring Functions}
We benchmark five post-hoc OOD scoring algorithms operating on internal representations of the frozen network, summarized in Table~\ref{tab:detector_formulations}.

\begin{table*}[!t]
\caption{Mathematical formulations, operating spaces, and hyperparameter specifications for evaluated post-hoc OOD detection algorithms.}
\label{tab:detector_formulations}
\centering
\footnotesize
\setlength{\tabcolsep}{4pt}
\resizebox{\columnwidth}{!}{%
\begin{tabular}{lllp{6.8cm}}
\toprule
\textbf{Detection Algorithm} & \textbf{Operating Representation} & \textbf{Mathematical Scoring Function $S(x)$} & \textbf{Implementation Details and Hyperparameters} \\
\midrule
\textbf{Maximum Softmax Probability (MSP) \cite{hendrycks2017baseline}} & Softmax Probabilities $p(x) \in [0, 1]^K$ & $S_{\text{MSP}}(x) = \max_{k \in \{1,\dots,K\}} \frac{\exp(z_k(x))}{\sum_{j=1}^K \exp(z_j(x))}$ & Baseline confidence scoring; bounded in $[1/K, 1]$. \\
\addlinespace[2pt]
\textbf{Energy-Based Detection \cite{liu2020energy}} & Raw Logits $z(x) \in \mathbb{R}^K$ & $S_{\text{Energy}}(x) = T \cdot \log \sum_{k=1}^K \exp(z_k(x) / T)$ & Post-hoc free energy score with $T = 1.0$; mathematically unbounded on $(-\infty, +\infty)$; not treated as a calibrated density estimator. \\
\addlinespace[2pt]
\textbf{ODIN Perturbation Variant \cite{liang2018enhancing}} & Input Space $\tilde{x} \in [-1, 1]^{H \times W \times C}$ & $S_{\text{ODIN}}(x) = \max_k \frac{\exp(z_k(\tilde{x}) / T)}{\sum_j \exp(z_j(\tilde{x}) / T)}$ & Evaluated variant uses sign gradient of maximum logit: $\tilde{x} = x + \varepsilon \operatorname{sign}(\nabla_x \max_k z_k(x))$ with $T = 1000.0, \varepsilon = 0.0014$. \\
\addlinespace[2pt]
\textbf{Ledoit--Wolf Mahalanobis \cite{lee2018simple}} & Latent Features $f(x) \in \mathbb{R}^{1280}$ & $S_{\text{Maha}}(x) = -\sqrt{\max(d_M^2(x), 0)}$ & Uses Ledoit--Wolf shrinkage covariance $\mathbf{\Sigma}$; square root is a strictly monotonic transformation preserving AUROC and ranking. \\
\addlinespace[2pt]
\textbf{Modified OpenMax Variant \cite{bendale2016towards}} & Latent Features $f(x) \in \mathbb{R}^{1280}$ & $S_{\text{OpenMax}}(x) = 1 - P(y = \text{unknown} \mid x)$ & Evaluated variant fits Weibull distributions to 1280-dim global features ($\alpha = 2, \eta = 20$) with rank weighting leaving top-1 logit unaltered. \\
\bottomrule
\end{tabular}%
}
\end{table*}

\subsection{Operational Energy--MSP Dual-Gate Cascade}
In production sorting lines, continuous scoring functions must trigger discrete accept/reject decisions. The Energy score is a post-hoc logit-based OOD scoring function motivated by an energy-based formulation; it is not treated here as a calibrated estimate of the true input marginal density. To combine the discriminative characteristics of logit-based Energy scoring with the confidence calibration of MSP, we evaluate an operational Boolean Dual-Gate cascade:
\begin{equation}
\operatorname{ACCEPT}(x) = \mathbb{I}(S_{\text{Energy}}(x) \ge \tau_E \;\land\; S_{\text{MSP}}(x) \ge \tau_M), \qquad \operatorname{REJECT}(x) = 1 - \operatorname{ACCEPT}(x).
\end{equation}
An exploratory linear-margin proxy, $S_{\text{proxy}}(x) = (S_{\text{Energy}}(x) - \tau_E) + (S_{\text{MSP}}(x) - \tau_M)$, facilitates continuous ROC analysis. Because raw Energy scores exhibit empirical variance ($\sigma \approx 2.15$) twelve times larger than that of bounded MSP ($\sigma \approx 0.18$), this proxy is scale-dependent and dominated by Energy variance (Spearman rank correlation $\rho > 0.992$). All operational screening metrics evaluate strictly the discrete Boolean conjunction rule.

\subsection{Validation Calibration and Evaluation Metrics}
Operating thresholds are locked prospectively on the independent primary validation partition ($N = 227$), targeting $95.15\%$ known retention ($11/227$ rejected): Standalone Energy $\tau_E = 1.054579$ ($s_{(11)}$); Dual-Gate Cascade $\tau_E = 0.918451 \land \tau_M = 0.582610$ ($s_{(7)}$ on both constituent gates, yielding $|R_E|=6, |R_M|=6, |R_E \cap R_M|=1 \implies |R_{\text{union}}|=11$). 

We report threshold-independent AUROC (OOD positive class, anomaly score $-S(x)$) and retrospective FPR@95 (OOD accepted at $\ge 95\%$ known retention). Under locked validation thresholds ($\tau_{95}$), we evaluate operational metrics: (1) \textbf{Unknown False Acceptance Rate (UFAR):} proportion of OOD images accepted into known distribution; (2) \textbf{UFAR-Fresh:} proportion of OOD images accepted and assigned one of the classifier's Fresh-class output labels ($\hat{y} \in \{0, 1, 2\}$), disclaiming commercial grade or biological freshness certification; and (3) \textbf{Known-Class Rejection Rate (KRR):} in-distribution abstention rate.

\subsection{Uncertainty Quantification Protocol}
Point estimates are accompanied by 95\% non-parametric bootstrap percentile intervals ($B = 1,000$ resamples with replacement, `\texttt{seed=42}'), stratified by class and condition groups. Crucially, Cohorts A--C share the FruitVision known reference partition ($N_{\text{id}} = 4,157$). Because bootstrap resampling was executed independently per cohort, marginal confidence intervals are not paired intervals for cohort differences, and these cohorts should not be treated as independent experimental replications.

\section{Experimental Results}
\label{sec:results}

\subsection{Closed-Set Classification Performance Across External Domains (RQ1)}
We evaluate the frozen MobileNetV2 classifier across the six known classes ($K = 6$) under laboratory conditions (Partition A) versus external repository conditions (Partition B), reported in Table~\ref{tab:closed_set_performance}.

\begin{table*}[!t]
\caption{Closed-set classification performance of frozen MobileNetV2 across primary laboratory and external photographic domains. Brackets indicate 95\% non-parametric bootstrap percentile intervals ($B = 1,000$, seed 42).}
\label{tab:closed_set_performance}
\centering
\footnotesize
\setlength{\tabcolsep}{4.5pt}
\resizebox{\columnwidth}{!}{%
\begin{tabular}{llcccc}
\toprule
\textbf{Evaluation Partition} & \textbf{Upstream Source Description} & \textbf{$N$} & \textbf{Top-1 Accuracy (\%)} & \textbf{Macro-Averaged F1 (\%)} & \textbf{Balanced Accuracy (\%)} \\
\midrule
\textbf{Partition A (ID Test)} & Kaggle Primary Benchmark (Lab) & 227 & 96.92 [94.71, 98.68] & 96.87 [94.65, 99.05] & 97.09 [94.80, 98.98] \\
\addlinespace[2pt]
\textbf{Partition B (Pooled)} & FruitVision ($N=4,157$) + Sultana ($N=1,193$) & 5,350 & 68.36 [67.20, 69.40] & 67.28 [66.21, 68.32] & 68.66 [67.55, 69.69] \\
\textbf{FruitVision B Subset} & FruitVision Known subset (mixed capture settings) & 4,157 & 69.14 [68.13, 70.17] & 66.75 [65.62, 68.01] & 69.56 [68.52, 70.58] \\
\textbf{Sultana B Subset} & Sultana Known subset (retail and field) & 1,193 & 65.63 [63.03, 68.15] & 65.77 [63.32, 68.13] & 65.55 [63.11, 67.99] \\
\bottomrule
\end{tabular}%
}
\end{table*}

Under laboratory capture, the model achieves 96.92\% top-1 accuracy (220/227 correct), confirming strong baseline extraction. On external images (Partition B), accuracy drops by 28.56 percentage points to 68.36\% (3,657/5,350 correct), with macro-F1 falling to 67.28\%. Complete non-overlap of 95\% bootstrap intervals reflects substantial descriptive separation. 

Class-wise inspection reveals an asymmetric collapse: the lowest accuracy occurs on Fresh Apple at 32.12\% (310/965 correct), systematically confused with Rotten Apple (393 / 40.73\%) and Rotten Orange (255 / 26.42\%). Qualitative inspection reveals natural surface striations, lenticel spots, and specularity characteristic of varietal russeting; we present this texture correlation strictly as an untested observational hypothesis, distinguishing qualitative visual inspection from causal evidence. In contrast, Rotten Banana achieves 99.40\% accuracy (826/831), and Fresh Orange achieves 87.53\% (828/946).

\subsection{Source-Stratified Out-of-Distribution Discrimination (RQ2)}
Table~\ref{tab:ood_benchmark} reports threshold-independent AUROC and retrospective FPR@95 across Cohorts A--D.

\begin{table*}[!t]
\caption{Source-stratified out-of-distribution detection benchmark across Cohorts A--D. AUROC (\%) and retrospective FPR@95 (\%) evaluate anomaly scoring functions under source-matched conditions. Brackets indicate 95\% bootstrap percentile intervals ($B = 1,000$, seed 42).}
\label{tab:ood_benchmark}
\centering
\footnotesize
\setlength{\tabcolsep}{3pt}
\resizebox{\columnwidth}{!}{%
\begin{tabular}{lcccccccc}
\toprule
& \multicolumn{2}{c}{\textbf{Cohort A: Semantic Shift}} & \multicolumn{2}{c}{\textbf{Cohort B: Condition Shift}} & \multicolumn{2}{c}{\textbf{Cohort C: Combined Shift}} & \multicolumn{2}{c}{\textbf{Cohort D: Semantic Shift}} \\
& \multicolumn{2}{c}{\textit{FruitVision B vs D1}} & \multicolumn{2}{c}{\textit{FruitVision B vs C}} & \multicolumn{2}{c}{\textit{FruitVision B vs D2}} & \multicolumn{2}{c}{\textit{Sultana B vs D1}} \\
\cmidrule(lr){2-3} \cmidrule(lr){4-5} \cmidrule(lr){6-7} \cmidrule(lr){8-9}
\textbf{Detection Algorithm} & \textbf{AUROC} & \textbf{FPR@95} & \textbf{AUROC} & \textbf{FPR@95} & \textbf{AUROC} & \textbf{FPR@95} & \textbf{AUROC} & \textbf{FPR@95} \\
\midrule
\textbf{Maximum Softmax Prob.\ (MSP)} & 72.42 & 82.74 & \textbf{58.77} & 92.85 & 76.47 & 76.40 & 59.77 & 94.30 \\
\textbf{Energy-Based Detection ($T=1.0$)} & \textbf{81.57} [80.81, 82.38] & \textbf{68.84} & 53.13 [52.16, 54.09] & 94.42 & \textbf{84.72} [83.85, 85.58] & \textbf{64.94} & 67.52 [65.67, 69.25] & 83.50 \\
\textbf{ODIN Perturbation Variant} & 72.30 & 87.43 & 56.73 & \textbf{91.13} & 78.11 & 83.85 & 60.02 & 94.25 \\
\textbf{Ledoit--Wolf Mahalanobis} & 81.41 & 79.00 & 47.91 & 98.07 & 80.08 & 78.49 & \textbf{68.15} & \textbf{81.20} \\
\textbf{Modified OpenMax Variant} & 72.46 & 91.28 & 55.35 & 95.20 & 77.22 & 88.12 & 64.38 & 94.75 \\
\midrule
\textbf{Dual-Gate Proxy (exploratory)} & 80.93 [80.13, 81.72] & 69.27 & 53.96 [52.96, 54.91] & 94.16 & 84.16 [83.28, 85.05] & 63.18 & 66.68 [64.98, 68.38] & 84.60 \\
\bottomrule
\end{tabular}%
}
\end{table*}

\noindent The empirical findings illuminate three key patterns:
\begin{enumerate}
\setlength{\itemsep}{0pt}\setlength{\parsep}{0pt}\setlength{\parskip}{0pt}\setlength{\topsep}{1pt}\setlength{\partopsep}{0pt}
\item \textbf{Comparative Performance on Semantic Shifts:} In Cohort A (FruitVision Semantic Shift), Energy scoring achieves the highest discrimination (AUROC = 81.57\% [80.81\%, 82.38\%], FPR@95 = 68.84\%), closely matched by Ledoit--Wolf Mahalanobis (81.41\%). Standard MSP (72.42\%) and ODIN (72.30\%) lag by ~9 percentage points. In Cohort C (Combined Shift), observed AUROC reaches 84.72\% (Energy) and 80.08\% (Mahalanobis); because Cohort C evaluates unseen species bearing dataset-reported treatment labels (Partition D2) rather than an isolated intervention, this reflects empirical cohort ordering and cannot isolate a causal treatment effect.
\item \textbf{Failure on Treatment-Labelled Shift (Cohort B):} When evaluated on known species bearing dataset-reported preservative labels (FruitVision C), all evaluated methods fail: Energy AUROC collapses to 53.13\% [52.16\%, 54.09\%], MSP to 58.77\%, and Mahalanobis drops to 47.91\% (FPR@95 = 98.07\%). Tested RGB methods do not reliably separate treatment-labelled images from natural fruit.
\item \textbf{Sultana Source-Matched Semantic-Shift Evaluation (Cohort D):} On Sultana Semantic Shift, discrimination is lower than on FruitVision: Mahalanobis achieves a higher AUROC point estimate (68.15\%) than Energy (67.52\% [65.67\%, 69.25\%]) and MSP (59.77\%), with FPR@95 exceeding 81\% to 95\%. Visual resemblance among spherical green commodities (guava, jujube) and known produce represents an untested observational hypothesis; morphological overlap was not directly quantified.
\end{enumerate}

\subsection{Operational Gating and Image-Level Decision Performance (RQ3)}
Table~\ref{tab:operational_screening} details operational decision outcomes under prospective, locked validation thresholds ($\tau_{95}$, targeting 95.15\% validation coverage).

\begin{table*}[!t]
\caption{Operational decision performance under prospective validation-calibrated thresholds ($\tau_{95}$). Standalone Energy uses $\tau_E = 1.054579$; Dual-Gate evaluates $S_{\text{Energy}} \ge 0.918451 \land S_{\text{MSP}} \ge 0.582610$. Brackets indicate 95\% bootstrap percentile intervals ($B = 1,000$, seed 42).}
\label{tab:operational_screening}
\centering
\footnotesize
\setlength{\tabcolsep}{3pt}
\resizebox{\columnwidth}{!}{%
\begin{tabular}{llcccc}
\toprule
\textbf{Evaluation Cohort} & \textbf{Operational Configuration} & \textbf{Operational UFAR (\%)} & \textbf{UFAR-Fresh (\%)} & \textbf{Operational KRR (\%)} & \textbf{Retained Knowns (\%)} \\
\midrule
\textbf{Cohort A: Semantic} & Standalone Energy ($\tau_E = 1.054579$) & 13.18 [12.07, 14.34] (366) & 7.64 [6.70, 8.65] (212) & 41.42 [40.34, 42.58] (1,722) & 58.58 (2,435) \\
{\scriptsize $N_{\text{id}}=4,157, N_{\text{ood}}=2,776$} & Dual-Gate Cascade ($\tau_E \land \tau_M$) & 16.07 [14.77, 17.40] (446) & 10.48 [9.37, 11.60] (291) & 41.98 [40.75, 43.06] (1,745) & 58.02 (2,412) \\
\addlinespace[2pt]
\textbf{Cohort B: Condition} & Standalone Energy ($\tau_E = 1.054579$) & 53.94 [52.16, 55.66] (1,034) & 44.03 [42.41, 45.59] (844) & 41.42 [40.34, 42.58] (1,722) & 58.58 (2,435) \\
{\scriptsize $N_{\text{id}}=4,157, N_{\text{ood}}=1,917$} & Dual-Gate Cascade ($\tau_E \land \tau_M$) & 52.63 [50.81, 54.25] (1,009) & 42.83 [41.26, 44.34] (821) & 41.98 [40.80, 43.23] (1,745) & 58.02 (2,412) \\
\addlinespace[2pt]
\textbf{Cohort C: Combined} & Standalone Energy ($\tau_E = 1.054579$) & 6.86 [5.52, 8.28] (82) & 4.94 [3.85, 6.19] (59) & 41.42 [40.34, 42.63] (1,722) & 58.58 (2,435) \\
{\scriptsize $N_{\text{id}}=4,157, N_{\text{ood}}=1,195$} & Dual-Gate Cascade ($\tau_E \land \tau_M$) & 9.54 [7.86, 11.30] (114) & 7.28 [5.94, 8.70] (87) & 41.98 [40.89, 43.20] (1,745) & 58.02 (2,412) \\
\addlinespace[2pt]
\textbf{Cohort D: Semantic} & Standalone Energy ($\tau_E = 1.054579$) & 34.15 [32.35, 35.90] (683) & 21.45 [19.85, 23.10] (429) & 41.74 [39.31, 44.26] (498) & 58.26 (695) \\
{\scriptsize $N_{\text{id}}=1,193, N_{\text{ood}}=2,000$} & Dual-Gate Cascade ($\tau_E \land \tau_M$) & 34.90 [33.00, 36.95] (698) & 22.65 [20.95, 24.30] (453) & 44.01 [41.41, 46.61] (525) & 55.99 (668) \\
\addlinespace[2pt]
\textbf{Partition B Pooled} & Standalone Energy ($\tau_E = 1.054579$) & --- & --- & 41.50 [40.21, 42.69] (2,220) & 58.50 (3,130) \\
{\scriptsize $N_{\text{id}}=5,350$} & Dual-Gate Cascade ($\tau_E \land \tau_M$) & --- & --- & 42.43 [41.10, 43.76] (2,270) & 57.57 [56.24, 58.90] (3,080) \\
\bottomrule
\end{tabular}%
}
\end{table*}

Operational screening demonstrates severe practical trade-offs:
\begin{itemize}
\setlength{\itemsep}{0pt}\setlength{\parsep}{0pt}\setlength{\parskip}{0pt}\setlength{\topsep}{1pt}\setlength{\partopsep}{0pt}
\item \textbf{In-Distribution Abstention Burden:} Fixed validation thresholds calibrated to accept 95.15\% of laboratory validation images suffer substantial over-rejection when transferred to external repositories. Dual-Gate incurs a Known-Class Rejection Rate (KRR) of 41.98\% on FruitVision B and 44.01\% on Sultana B, pooling at 42.43\% [41.10\%, 43.76\%] (2,270/5,350 known images rejected). Retaining automated decisions on only 57.57\% of produce imposes a substantial operational abstention burden.
\item \textbf{Treatment Routing Failure:} In Cohort B, Dual-Gate accepts 42.83\% [41.26\%, 44.34\%] of treatment-labelled image records (821/1,917) and routes them into the classifier's Fresh-class output labels (Standalone Energy UFAR-Fresh is 44.03\%, 844/1,917). Under the evaluated operating thresholds, a substantial fraction of treatment-labelled image records were accepted and assigned one of the classifier's Fresh-class output labels; this outcome does not establish that their confidence-score distributions match those of untreated fresh fruit.
\item \textbf{Cross-Source Semantic Leakage:} In Cohort D, operational false acceptance rises to 34.90\% (Dual-Gate) with 22.65\% UFAR-Fresh (453/2,000 accepted as fresh knowns). Under locked thresholds, the evaluated score-based gate did not reliably reject all unseen-species records in Cohort D; score-distribution and feature-space overlap were not independently quantified.
\end{itemize}

\subsection{Perceptual Hash Near-Duplicate Sensitivity Analysis}
Auditing detector performance across deduplication tiers confirms stability: across all four cohorts and five detectors, the maximum observed change in AUROC between Primary ($N=13,465$) and Tier 2 ($N=13,382$) is 0.22 percentage points (Cohort C Mahalanobis: 80.08\% $\to$ 79.86\%). In Cohort A, Energy AUROC shifts marginally from 81.57\% (Primary) to 81.55\% (Tier 1) and 81.46\% (Tier 2); in Cohort B, it shifts from 53.13\% to 53.16\% and 53.19\%. In Cohort D, metrics remain completely invariant. Relative detector rankings remain strictly invariant across all evaluation tiers.

\section{Discussion and Limitations}
\label{sec:discussion_limitations}

\subsection{Domain Collapse and Weak Discrimination on Treatment-Labelled Images}
Our empirical findings reveal fundamental failure modes in deploying closed-world vision models to post-harvest sorting. The 28.56 percentage-point degradation in top-1 accuracy (from 96.92\% to 68.36\%) demonstrates the fragility of deep convolutional backbones to unconstrained environmental shifts. In curated laboratory benchmarks, models often exploit spurious background correlations (e.g., uniform illumination, neutral mats). When deployed across external markets, mobile camera optics, ambient shadowing, and surface specularities degrade feature extraction. The severe collapse on Fresh Apples (32.12\% accuracy; 393 misclassified as Rotten Apple, 255 as Rotten Orange) coincides with visual surface variations; we present the correlation between natural surface striations, lenticel spots, and decay misclassifications strictly as an untested observational hypothesis, distinguishing qualitative visual inspection from causal evidence.

Furthermore, the breakdown of OOD detection on dataset-reported treatment-labelled produce in Cohort B (all evaluated detectors AUROC $\le 58.77\%$; 42.83\% Dual-Gate UFAR-Fresh) highlights the challenge of identifying chemical treatment status from standard RGB imagery without independent biochemical assays. On these dataset-assigned labels, the frozen convolutional backbone and post-hoc scoring functions do not achieve reliable separation from natural knowns (e.g., Mahalanobis AUROC = 47.91\%, Energy = 53.13\%). Rather than asserting general visual indistinguishability, we report strictly that the evaluated post-hoc RGB detectors fail on this repository-assigned treatment cohort. Future work could evaluate whether complementary sensing modalities improve screening under these dataset-reported condition shifts, such as: (1) near-infrared and hyperspectral spectroscopy to assess internal tissue properties, (2) volatile organic compound (VOC) headspace gas sensing to detect chemical markers \cite{sattar2024sensornet}, and (3) ambient storage and cold-chain telemetry.

\subsection{Operational Trade-offs in Decision Gating}
Our benchmark demonstrates a critical divergence between continuous ranking metrics and offline image-level acceptance and rejection outcomes. While Energy scoring achieves 81.57\% AUROC on semantic shift, prospective gating under locked validation thresholds incurs an operational Known-Class Rejection Rate (KRR) of 42.43\% across external known image records (2,270/5,350 image records rejected). If these image-level rejection rates transferred to an operational line, they could increase manual-review demand; this study did not measure commodity-level diversion, throughput, labour, or grading outcomes. While the Dual-Gate cascade suppresses false acceptance on semantic shifts, both constituent gates operate downstream of the same MobileNetV2 features and remain fundamentally bounded by visual representation capacity.

\subsection{Limitations and Threats to Validity}
To preserve research transparency, we articulate key experimental boundaries:
\begin{enumerate}
\setlength{\itemsep}{0pt}\setlength{\parsep}{0pt}\setlength{\parskip}{0pt}\setlength{\topsep}{1pt}\setlength{\partopsep}{0pt}
\item \textbf{Image Records vs.\ Specimen Independence:} Evaluated datasets lack biological fruit identifiers. Multiple records may represent rotations or burst captures of shared specimens. Although pHash sensitivity auditing across Tier 1 (excluding 56 cross-partition collisions at $r \le 8$) and Tier 2 (excluding an additional 27 intra-partition multi-label collisions) confirmed detector ranking invariance, mathematical hashing cannot prove biological independence. All bootstrap summaries are conditional on image-level resampling exchangeability.
\item \textbf{Dataset-Reported Treatment Labels without Independent Chemical Assays:} Condition labels for Cohort B and C rely on FruitVision metadata \cite{bijoy2025fruitvision}. No independent gas chromatography--mass spectrometry (GC-MS) or tissue assays were conducted by the authors. We report strictly that evaluated post-hoc vision gates exhibited weak separation on this dataset-reported cohort.
\item \textbf{Residual Within-Repository Confounding:} While source matching reduces cross-repository hardware divergence, within-repository variations (different market stalls, dates, mobile phones) remain possible.
\item \textbf{Statistical Conditioning across Cohorts:} Cohorts A--C share the FruitVision Known reference partition ($N_{\text{id}} = 4,157$). Because bootstrap resampling was executed independently per cohort, marginal confidence intervals are not paired intervals for cohort differences, and cohorts are not independent replications.
\item \textbf{Offline Evaluation without Physical Conveyance:} This study benchmarks static digital images without evaluating high-speed conveyor transport, motion blur, variable camera trigger latency, or pneumatic sorting physics.
\item \textbf{Algorithmic Implementation Variants:} As detailed in Section~\ref{sec:methodology}, the evaluated ODIN implementation optimizes the sign gradient of the maximum logit, and OpenMax fits Weibull tails to 1280-dimensional global features. Results reflect these verified implementations.
\item \textbf{Exploratory Margin Proxy Scale Dependence:} The continuous Energy--MSP linear-margin proxy (exploratory) is dominated by Energy score variance ($\sigma \approx 2.15$ vs.\ $0.18$, $\rho > 0.992$); operational decision metrics rely strictly on the scale-immune Boolean conjunction rule.
\end{enumerate}

\section{Conclusion}
\label{sec:conclusion}

This study conducted a source-stratified empirical benchmark evaluating closed-set fruit-quality classification and post-hoc out-of-distribution detection across $N = 13,465$ photographic records under domain, semantic species, and dataset-reported treatment shifts. Our findings provide empirical answers to the three core research questions:
\begin{enumerate}
\setlength{\itemsep}{0pt}\setlength{\parsep}{0pt}\setlength{\parskip}{0pt}\setlength{\topsep}{1pt}\setlength{\partopsep}{0pt}
\item \textbf{Resolution of RQ1 (External Domain Robustness):} Evaluating a frozen MobileNetV2 classifier on external repository imagery reveals a marked 28.56 percentage-point degradation in top-1 accuracy (from 96.92\% [94.71\%, 98.68\%] on laboratory test data to 68.36\% [67.20\%, 69.40\%] on external knowns). This degradation is accompanied by a pronounced class-wise failure on Fresh Apple (32.12\% accuracy; 310/965 correct), where natural surface striations, lenticel spots, and specularity are observed to coincide with decay misclassifications (393 Rotten Apple, 255 Rotten Orange), which we present strictly as an untested observational hypothesis.
\item \textbf{Resolution of RQ2 (Comparative Out-of-Distribution Discrimination):} Under source-stratified semantic species shifts, Energy scoring (AUROC = 81.57\% on FruitVision, 67.52\% on Sultana) and Ledoit--Wolf Mahalanobis distance (81.41\% and 68.15\%) achieve higher AUROC point estimates than baseline Maximum Softmax Probability (72.42\% and 59.77\%). However, on dataset-reported chemical preservative shift (Cohort B), all evaluated post-hoc scoring algorithms exhibit weak or near-chance discrimination (Energy AUROC = 53.13\%, Mahalanobis = 47.91\%, FPR@95 $\ge 91.13\%$). Post-hoc scoring methods on standard RGB inputs do not reliably separate treatment-labelled images from natural produce.
\item \textbf{Resolution of RQ3 (Operational Feasibility and Rejection Trade-offs):} Prospective validation-calibrated gating ($\tau_{95}$) achieves partial screening on pure semantic species shift in Cohort A (Dual-Gate UFAR = 16.07\%) and on combined species-and-treatment shift in Cohort C (Dual-Gate UFAR = 9.54\%), but incurs a substantial in-distribution abstention penalty, rejecting 42.43\% [41.10\%, 43.76\%] of external known-class image records (2,270/5,350) and retaining automated decisions on only 57.57\% of image records. Crucially, 42.83\% [41.26\%, 44.34\%] of dataset-reported treatment-labelled image records (821/1,917) are accepted into the classifier's Fresh-class output labels.
\end{enumerate}

\noindent This investigation demonstrates the critical necessity of source-stratified evaluation in agricultural AI and exposes the visual failure boundaries of standalone RGB computer vision. Future work could evaluate whether multi-modal architectures integrating optical imaging with internal spectroscopy or chemical headspace sensing improve screening under these condition shifts.


\section*{Acknowledgements}
The authors acknowledge the creators and curators of the Kaggle Fruit Quality, FruitVision, and Sultana datasets for making their photographic archives available for academic research.

\bibliographystyle{elsarticle-num}
\bibliography{bibliography/references}

\end{document}
